\chapter{Dataset Description}

Supervised detection needs labelled examples of both normal and malicious runtime
behaviour in exactly the format the sensor produces. No public dataset provides one
row per kernel event with Kubernetes context, behaviour-window counters and an attack
category (Section~2.5), so three sources were combined:
\begin{enumerate}
    \item a \textbf{synthetic runtime-telemetry dataset} generated for this project, which
          provides all labelled attacks of the nine categories;
    \item the public \textbf{BETH} dataset of real eBPF kernel events from cloud honeypots;
    \item a real \textbf{Kubernetes container-escape} recording made with Tetragon
          (Splunk Attack Data, Cisco Isovalent).
\end{enumerate}
All sources were converted to the same schema and merged into one final training set
and one final test set (Section~4.5).

\section{Event Schema}

Every row describes one kernel event after Kubernetes enrichment, plus the behaviour of
the same container in the preceding 60 seconds. Table~\ref{tab:schema} lists the
columns. The same column names are produced by the live agent, so a model trained on
these rows can be applied to live events without any conversion.

\begin{longtable}{>{\ttfamily\footnotesize}p{4.0cm} >{\raggedright\arraybackslash}p{9.1cm}}
\caption{Columns of the runtime-telemetry dataset}\label{tab:schema}\\
\toprule
\normalfont\textbf{Column} & \textbf{Description} \\
\midrule
\endfirsthead
\toprule
\normalfont\textbf{Column} & \textbf{Description} \\
\midrule
\endhead
event\_type & exec, open, connect, ptrace, setuid, mount, unshare, chmod, module, memfd, setns \\
comm, parent\_comm & Process name and parent process name (15 characters, as in the kernel) \\
uid & Numeric user ID of the process \\
path & Executed binary, opened file, UNIX socket, mount target or memfd name \\
cmdline, argc & Command-line arguments (max.\ 8 $\times$ 127 bytes) and their count \\
flags, arg2 & open flags, chmod mode, unshare flags, ptrace request or target uid \\
dport, daddr & Destination port and address of a connection \\
namespace & Kubernetes namespace of the pod \\
win\_exec\_count & Executions in the container in the last 60 s \\
win\_distinct\_bins & Distinct binaries executed in the window \\
win\_connect\_count & Outgoing connections in the window \\
win\_distinct\_dst & Distinct destinations (address:port) in the window \\
win\_sensitive\_opens & Opens of sensitive files in the window \\
win\_shell\_count & Shell executions in the window \\
win\_new\_binary & 1 if the binary was never seen in this workload after a 5-minute warm-up \\
is\_malicious & Ground truth: 0 = benign, 1 = malicious \\
attack\_category & benign or one of nine attack categories \\
mitre\_technique & MITRE ATT\&CK technique ID of the template \\
template & Name of the generating template / scenario \\
is\_edge\_case & 1 for benign-but-suspicious or malicious-but-benign-looking rows \\
\bottomrule
\end{longtable}

\section{Synthetic Runtime-Telemetry Dataset}

The generator (\texttt{ml/generate\_dataset.py}) produces events from \textbf{52 template
families} (29 benign, 23 attack) with about 300 concrete command and event variants.
IP addresses, ports, file names, user IDs, namespaces and behaviour windows are
randomised, so no two rows are identical by construction. The 17 command patterns of
the original v0 dataset are kept as templates. The command lines pass through the same
truncation as the kernel probe (eight arguments of 127 bytes).

\subsection{Composition}

The training set was generated with 30,000 benign and 10,000 malicious events (seed 7);
an independent test set with 12,000 benign and 8,000 malicious events uses a different
random seed (2026). Table~\ref{tab:categories} shows the distribution of the training
set over categories and the MITRE techniques covered.

\begin{table}[H]
\centering
\caption{Synthetic training data by category (seed 7, 40,000 rows)}
\label{tab:categories}
\small
\begin{tabular}{lrrl}
\toprule
\textbf{Category} & \textbf{Templates} & \textbf{Rows} & \textbf{Main MITRE techniques} \\
\midrule
Benign                & 29 & 30,000 & -- \\
Reverse shell         & 5  & 2,135  & T1059.004, T1571 \\
Dropper / payload     & 3  & 1,432  & T1105, T1204.002 \\
Discovery             & 3  & 1,178  & T1613, T1046, T1190 \\
Privilege escalation  & 3  & 1,126  & T1068, T1548.001, T1098 \\
Container escape      & 2  & 1,033  & T1611, T1547.006 \\
Credential access     & 2  & 1,021  & T1528, T1003.008, T1552 \\
Crypto-miner          & 2  & 865    & T1496 \\
Defense evasion       & 2  & 785    & T1070, T1620 \\
Lateral movement      & 1  & 425    & T1210, T1021 \\
\midrule
\textbf{Total}        & \textbf{52} & \textbf{40,000} & \\
\bottomrule
\end{tabular}
\end{table}

The event types in the training set are dominated by process executions (23,262),
followed by connections (6,889), sensitive file opens (5,162), chmod (1,256), memfd
(1,129), setuid (1,017), mount (455), ptrace (385), setns (205), unshare (162) and
kernel-module loads (78).

\subsection{Edge Cases}

More than half of the rows (24,589) are deliberate edge cases, flagged with
\texttt{is\_edge\_case = 1}, so the model learns the difficult boundary between benign
and malicious behaviour instead of trivial keywords:
\begin{itemize}
    \item \textbf{Benign but suspicious-looking:} \texttt{runc} and container runtimes
          calling mount and namespace system calls, \texttt{nc -z} and \texttt{/dev/tcp}
          wait-for loops, health probes with \texttt{curl \ldots{} || exit 1}, base64-decoded
          configuration, Go and C builds executing from \texttt{/tmp}, CoreDNS reading its
          service-account token, the API server reading \texttt{/etc/kubernetes/pki},
          \texttt{sshd} reading \texttt{/etc/shadow}, PostgreSQL writing \texttt{/dev/shm},
          CSI drivers mounting block devices, \texttt{strace} in development pods and
          \texttt{kubectl} or \texttt{helm} in CI namespaces.
    \item \textbf{Malicious but benign-looking:} reverse shells on ports 443, 80 and 53,
          base64-obfuscated payloads, miners renamed to \texttt{nginx} or
          \texttt{kworkerds}, fileless execution via \texttt{memfd\_create},
          \texttt{curl \ldots{} | bash} from GitHub, slow-and-low connections and
          reconnaissance spawned by a web server.
\end{itemize}

\section{BETH Dataset}

BETH (BPF-extended tracking honeypot)~\cite{highnam2021} contains about eight million
kernel events recorded with eBPF on 23 honeypot hosts of a major cloud provider. Each
event records the timestamp, process and parent IDs, user ID, mount namespace, process
name, host name, system-call name, arguments and two labels: \texttt{sus} (suspicious)
and \texttt{evil}. The attacks are a real botnet compromise: an SSH login, system
reconnaissance (\texttt{uname}, \texttt{/proc/cpuinfo}, \texttt{free}), download of a
payload from a raw IP address with \texttt{wget}, preparation and launch of a
crypto-miner from a hidden \texttt{/tmp/.X25-unix} directory and a mass SSH scan of
thousands of hosts.

\subsection{Conversion to the Project Schema}

A converter (\texttt{ml/convert\_external.py}) keeps only the system calls that our
sensor also captures and maps them to the project schema:
\begin{itemize}
    \item kernel calls are mapped as shown in Table~\ref{tab:bethmap};
    \item behaviour-window columns are recomputed by replaying the events through the
          agent's own window logic, keyed by host and mount namespace;
    \item \texttt{is\_malicious} = \texttt{evil}; BETH has no Kubernetes namespace and no
          attack categories.
\end{itemize}
\begin{table}[H]
\centering
\caption{Mapping of BETH system calls to the project's event types}
\label{tab:bethmap}
\small
\begin{tabular}{>{\ttfamily}l l >{\raggedright\arraybackslash}p{6.2cm}}
\toprule
\normalfont\textbf{BETH event} & \textbf{Event type} & \textbf{Fields used} \\
\midrule
execve & exec & path, argv; the calling process becomes the parent \\
security\_file\_open & open & path, O\_* flags \\
connect & connect & port and address, or UNIX socket path \\
setuid, setreuid & setuid & target user ID \\
chmod, fchmod & chmod & path, mode \\
mount & mount & source, target, filesystem type \\
memfd\_create & memfd & memfd name \\
\bottomrule
\end{tabular}
\end{table}

Of 3,607,974 raw events in eight files, 938,432 events map to the sensor's event types
(147,629 malicious). The conversion runs in parallel on all CPU cores and takes about
two minutes.

\section{Kubernetes Container-Escape Recording}

The Splunk Attack Data repository~\cite{splunkattack} provides Tetragon
\texttt{process\_exec} events recorded during a container escape in a Kubernetes cluster:
an attacker opens a shell in a privileged pod, runs
\texttt{nsenter --target 1} (entering all namespaces of the host's PID~1), obtains a shell on the
host, mounts \texttt{/proc}, reads \texttt{/root/.ssh/authorized\_keys} and performs
reconnaissance. The 25 events contain the binary, arguments, parent, user ID and Pod
metadata. Each event was labelled by hand: 7 attacker actions are malicious; 15
login-shell profile commands that run automatically after the escape and 3 pod start-up
events are benign.

\section{Final Training and Test Sets}

The final split (\texttt{ml/build\_final\_dataset.py}) follows two principles:
(1) the BETH data is split \emph{by host}, never by row, so near-identical events of one
machine cannot appear in both sets; (2) no test data is used for training or tuning.
Only benign BETH events of the development hosts (\texttt{evil = 0} and
\texttt{sus = 0}, duplicates removed) are added to training, because the few malicious
events on those hosts are mostly session noise and BETH has no attack categories.

\begin{table}[H]
\centering
\caption{Final training and test datasets}
\label{tab:finalsplit}
\small
\begin{tabular}{llrrr}
\toprule
\textbf{Split} & \textbf{Source} & \textbf{Rows} & \textbf{Malicious} & \textbf{Benign} \\
\midrule
\multirow{2}{*}{train.csv} & Synthetic (seed 7) & 40,000 & 10,000 & 30,000 \\
 & BETH benign, 15 development hosts & 25,000 & 0 & 25,000 \\
\cmidrule{2-5}
 & \textbf{Total} & \textbf{65,000} & \textbf{10,000} & \textbf{55,000} \\
\midrule
\multirow{3}{*}{test.csv} & Synthetic (seed 2026) & 20,000 & 8,000 & 12,000 \\
 & BETH held-out hosts 217 and 105 & 262,825 & 147,308 & 115,517 \\
 & Kubernetes escape (Tetragon) & 25 & 7 & 18 \\
\cmidrule{2-5}
 & \textbf{Total} & \textbf{282,850} & \textbf{155,315} & \textbf{127,535} \\
\bottomrule
\end{tabular}
\end{table}

\textbf{Characteristics to keep in mind when reading the results:}
\begin{itemize}
    \item All attack examples used for training are synthetic; the real data contributes
          only benign behaviour.
    \item 98\% of the malicious BETH test events (144,801 of 147,308) belong to one
          behaviour: the SSH-scanning bot. Results on BETH therefore mainly measure whether
          this scanner is recognised.
    \item BETH labels every event of the attacker's login session as evil, including
          routine commands, which lowers the achievable per-event recall.
    \item About 33\% of the test feature vectors also occur in the training data. We
          verified that 91,820 of these 93,946 rows are benign file opens on BETH hosts that
          run the same operating-system image, not leaked attacks.
\end{itemize}
